\documentclass[11pt]{article}

\usepackage[margin=1in]{geometry}
\usepackage{amsmath, amssymb}
\usepackage{enumitem}
\usepackage{actuarialangle}
\usepackage{setspace}

\setlength{\parindent}{0pt}
\setlength{\parskip}{8pt}

\begin{document}

\begin{center}
    {\Large \textbf{Annuities}}\\
\end{center}

\vspace{10pt}

\hrule

\medskip

\textbf{Problem 1.} An annuity-immediate provides annual payments of $10$ for $20$ years. Immediately following the $11^{\text{th}}$ payment, the annuity is exchanged for a perpetuity-immediate of equal value with semiannual payments. The present values at the time of exchange are based on an annual effective interest rate of $6\%$.

The first payment of the perpetuity is $K$, and each subsequent payment is $0.5\%$ larger than the previous payment.

Calculate $K$.

\medskip

\hrule


\medskip


\textbf{Problem 2.} A loan of $\$60{,}000$ is repaid over ten years with semiannual payments at a nominal interest rate of $6\%$ compounded semiannually.

The first payment is $\$3{,}000$ and is paid 6 months from the date of the loan. For the first five years, each subsequent payment will be $\$300$ more than the previous payment. After the $10^{\text{th}}$ payment, each succeeding semiannual payment will be reduced by $X$ until the loan is completely paid.

Calculate $X$.

\medskip

\hrule


% ------------------ Problem ------------------
\textbf{Problem 3.} A series of payments of $1$ is made at times
\[
t=10,14,18,\ldots,34.
\]
Let $i$ denote the annual effective interest. Find an expression od the present value at time $0$ of this payment stream as a function of $i$.

\begin{enumerate}[label=(\Alph*), leftmargin=2.2em]
    \item $\displaystyle \frac{a_{\actuarialangle{34}i}-a_{\actuarialangle{6}i}}{s_{\actuarialangle{4}i}}$
    \item $\displaystyle \frac{a_{\actuarialangle{35}i}-a_{\actuarialangle{7}i}}{s_{\actuarialangle{4}i}}$
    \item $\displaystyle \frac{\ddot{a}_{\actuarialangle{34}i}-\ddot{a}_{\actuarialangle{6}i}}{s_{\actuarialangle{4}i}}$
    \item $\displaystyle \frac{a_{\actuarialangle{34}i}-a_{\actuarialangle{6}i}}{\ddot{s}_{\actuarialangle{4}i}}$
    \item $\displaystyle \frac{\ddot{a}_{\actuarialangle{34}i}-\ddot{a}_{\actuarialangle{6}i}}{\ddot{s}_{\actuarialangle{4}i}}$
\end{enumerate}

\medskip
\hrule






\newpage

\hrule

\medskip

\textbf{Problem 4.} A perpetuity makes continuous payments. At time $t$, the perpetuity pays at an annual rate of $3t$, where time is measured in years. 

At am annual effective interest rate of $5\%$, calculate the present value of the perpetuity.

\medskip

\hrule

\medskip

\textbf{Problem 5.} Jen invests $\$X$ at the beginning of each year for $6$ years at an annual effective interest rate of $5\%$. The interest credited at the end of each year is reinvested at an annual effective interest rate of $7\%$.

Bill invests $\$X$ at the end of each year for $10$ years at an annual effective interest rate of $i\%$. The accumulated value of Bill's investments at the end of $10$ years is twice as much as the accumulated value of Jen's investments at the end of $6$ years.

Calculate $i$.

\medskip
\hrule


\medskip

\newpage

\textbf{Problem A.} \$1{,}000 is deposited into a savings account at time $t=0$. No other amounts are deposited. The accumulated amount in the account at any time is given by
\[
A(t) = 1000\left(1 + \frac{2t}{35}\right)^2.
\]

At what time $t$ is the force of interest equal to $0.04$?

\medskip

\hrule

\medskip

\textbf{Problem B.} Joe deposits \$10 today and another \$30 in five years into a fund paying simple interest of $11\%$ per year.

Tina will make the same two deposits, but the \$10 will be deposited $n$ years from today and the \$30 will be deposited $2n$ years from today. Tina's deposits earn an annual effective interest rate of $9.15\%$.

At the end of $10$ years, the accumulated amount of Tina's deposits equals the accumulated amount of Joe's deposits.

Calculate $n$.

\medskip

\hrule


\medskip


\textbf{Problem C.} Jim deposits \$100 into an account that credits:
\begin{itemize}[leftmargin=2em]
    \item An effective rate of discount of $3\%$ for the first year;
    \item An effective rate of discount of $6\%$ for the second year;
    \item An effective rate of discount of $9\%$ for the third year; and
    \item $\delta_t = 0.02t$ for the fourth and fifth years.
\end{itemize}

How much is in the account at the end of the fifth year?

\medskip

\hrule





\end{document}
